\chapter{Literature Review}
\label{ch:literature-review}

\lipsum[11-12]


\section{Background}
\label{sec:background}

\lipsum[13]

\section{why?}
\label{sec:why}

\lipsum[13]
\subsection{sub why?}
\label{sec:sub-why}

\lipsum[13]

\subsection{sub why?}
\label{sec:sub-why28}

\lipsum[13]
\subsection{sub why?}
\label{sec:sub-why27}

\lipsum[13]

\subsection{sub why?}
\label{sec:sub-why26}

\lipsum[13]

\subsection{sub why?}
\label{sec:sub-why25}

\lipsum[13]

\subsection{sub why?}
\label{sec:sub-why24}

\lipsum[13]

\subsection{sub why?}
\label{sec:sub-why23}

\lipsum[13]

\subsection{sub why?}
\label{sec:sub-why22}

\lipsum[13]

\subsection{sub why?}
\label{sec:sub-why21}

\lipsum[13]

\insertimage{architecture.png}{Sample system overview diagram}{fig:architecture}{1\textwidth}

As shown in Figure~\ref{fig:architecture}, the system consists of multiple
interconnected components. \lipsum[14]


\section{Related Work}
\label{sec:related-work}

\lipsum[15]

\inserttable{Comparison of related systems}{tab:related-work}{lccc}{
  \toprule
  \textbf{System} & \textbf{Year} & \textbf{Technology} & \textbf{Open Source} \\
  \midrule
  System Alpha & 2022 & Python / Flask      & Yes \\
  System Beta  & 2023 & Java / Spring Boot  & No  \\
  System Gamma & 2024 & Node.js / Express   & Yes \\
  \bottomrule
}

Table~\ref{tab:related-work} summarizes the key characteristics of existing
systems in the domain. \lipsum[16]


\section{Summary}
\label{sec:lit-summary}

\lipsum[17]

Several studies have contributed \abbr{ide} significantly to this field, including the
works by AuthorA et al.\ \insertref{knuth1984} and AuthorB et al.\
\insertref{lamport1994}. These references \abbrlong{html} provide the foundational \abbrfull{api} understanding
that this project \abbrshort{api} builds upon.
